\documentclass[10pt,a4paper]{amsart}
\usepackage[margin=19mm]{geometry}
\usepackage{amsmath,amssymb,mathtools}
\usepackage[scr=rsfs,cal=euler]{mathalfa}
\usepackage{xfrac}
\usepackage[unicode,hidelinks,bookmarks=false]{hyperref}
\newcommand{\ie}{i.e.\ }
\newcommand{\cf}{cf.\ }
\newcommand{\resp}{respectively\ }
\newcommand{\Subsection}{\subsection}
\providecommand{\nobreakdash}{\nobreak\hbox{-}}
\setlength{\emergencystretch}{2em}
\multlinegap0pt
\usepackage{amssymb}
\usepackage[normalem]{ulem}
\newtheorem{prop}{Proposition}
\newcommand{\N}{\mathbb{N}}
\title{Fixed Point Theorems for Set-Valued Maps with Contractive Orbits}
\author{Detelina Kamburova}
\date{Source: arXiv:2601.17380v3 (CC BY 4.0)}
\begin{document}
\maketitle

\noindent\emph{Source and licence.} Fixed Point Theorems for Set-Valued Maps with Contractive Orbits. Version arXiv:2601.17380v3 (CC BY 4.0), distributed by arXiv under the Creative Commons Attribution 4.0 International licence, http://creativecommons.org/licenses/by/4.0/, as recorded in the arXiv licence metadata for this exact version and shown on the abstract page https://arxiv.org/abs/2601.17380v3. Full licence notice: \texttt{licenses/arxiv-2601.17380-CC-BY-4.0.txt}.\par\smallskip

\noindent\textit{Editorial scope.} Counted unit: the article's prop t=>p (source0.tex lines 663--664). The article's complete original proof of it is delivered with it (source0.tex lines 665--680, one \texttt{proof} environment of 16 lines). No mathematics is written, repaired or re-derived: the statement and the proof are the article's own lines, in the article's own notation, and no retained line is substituted or re-flowed.  No further source-local context is needed: every cross-reference of every retained line resolves inside the two delivered spans.  Nothing was omitted from any retained span, so the package carries no omission record.  No retained line contains a citation, so the wrapper carries no bibliography and no bibliography entry.  No retained line contains a figure environment, an image-inclusion command or drawing code, so no image is delivered and the package declares no asset dependency.  Editorial formatting only, in addition: an AMS \texttt{amsart} preamble that restates the article's own declarations and macros used by the retained lines, namely the environments \texttt{prop} on separate counters, together with the standard packages those lines need.  The retained environments print in this document as Proposition 1 in order of appearance, whereas the article prints the same items with the numbering of its own full document.  The complete native source of the article as distributed in the arXiv e-print for this exact version is bundled unchanged as \texttt{sources/193311/source0.tex}.
\begin{prop} \label{t=>p}  Let $(X, \tau, h)$ be a Hausdorff premetric space. Let $S:X \rightrightarrows X$ be a set-valued map. Let $M:=\{x_i\}_{i=1}^{\infty}$ be an infinite f-t-contractive $S$-orbit. Then $M$ is a h-contractive $S$-orbit.
\end{prop}

\begin{proof} Since $\bar{x}$ is a strong accumulation point of $M$ then there exists a subsequence $\{x_{i_k}\}_{k=1}^{\infty}$ and whenever $a_k \in S(x_{i_k})$  $a_k \rightarrow \bar{x}$, i.e., the sequence
\[x_{i_1},a_1,x_{i_2},a_2,\dots,x_{i_k},a_k,\dots\]
converges to $\bar{x}$.
What remains to be proved is $h_{{S(x_{i_k}})}(x_{i_k}) \rightarrow 0$ as $k \rightarrow \infty$.

Fix $\varepsilon>0$. For all $k$ there exists $a_k^{\varepsilon} \in S(x_{i_k})$ such that
\begin{equation} \label{sup_eps}
h_{S(x_{i_k})}(x_{i_k})<h(a_k^{\varepsilon},x_{i_k})+\varepsilon.
\end{equation}

From (P3') it follows that $h(a_k^{\varepsilon}, x_{i_k}) \rightarrow 0$ as $k \rightarrow \infty$. Pick $K \in \N$ such that for all $k \geq K$ it holds that $h(a_k^{\varepsilon}, x_{i_k}) < \varepsilon$. From \ref{sup_eps} it follows that 
\[
h_{S(x_{i_k})}(x_{i_k})<2\varepsilon. 
\]
Since we chose $\varepsilon$ arbitrarily then $h_{S(x_{i_k})}(x_{i_k}) \rightarrow 0$ as $k \rightarrow \infty$. Therefore, $\bar{x}$ is a p-contractive accumulation point of $M$.
\end{proof}

\end{document}
